% One-page, questions-only practice exam based on 10162025155035.pdf
% (2025 Exam 2) and the topics in exam2mats.pdf.
% Original question designs, informed by review problems 3--4, 15--16,
% 26, 32, 39--40, and 43--50. No Green's theorem is needed.
\documentclass[12pt,letterpaper]{article}
\usepackage[utf8]{inputenc}
\usepackage[T1]{fontenc}
\usepackage[margin=0.8in]{geometry}
\usepackage{amsmath,amssymb}
\usepackage{newtxtext,newtxmath}
\usepackage{enumitem}

\setlength{\parindent}{0pt}
\setlength{\parskip}{0pt}
\pagestyle{empty}
\newcommand{\dd}{\mathop{}\!\mathrm{d}}

\begin{document}

\begin{enumerate}[label=\arabic*.,leftmargin=2em,labelsep=0.6em,
                  itemsep=1.2em,parsep=0pt,topsep=0pt]

\item Evaluate the following integral:
\[
  \int_0^1\int_{\sqrt{y}}^1 e^{x^3}\,\dd x\,\dd y.
\]

\item Suppose a solid lies inside the sphere $x^2+y^2+z^2=4$ and above the
cone $z=\sqrt{x^2+y^2}$. Its density is
$\delta(x,y,z)=z^2$. Set up, but do not evaluate, a triple integral
representing the mass of this solid in each coordinate system.

\begin{enumerate}[label=(\alph*),leftmargin=1.8em,topsep=0.4em,itemsep=0.2em]
  \item Cylindrical coordinates.
  \item Spherical coordinates, with $\varphi$ measured from the positive
        $z$-axis.
\end{enumerate}

\item Suppose a solid lies above the paraboloid $z=x^2+y^2$ and below the
paraboloid $z=8-x^2-y^2$. Set up a triple integral in an appropriate
coordinate system representing the volume of this solid, and evaluate it.

\item A wire follows the curve parametrized by
\[
  \vec r(t)=(t^2,3t,e^t),\qquad 0\leq t\leq1,
\]
and has density $\delta(x,y,z)=y+z$. Set up, but do not evaluate,
a line integral representing the mass of the wire.

\item Compute the work done by the force field $\vec F(x,y)=(-y,x)$ on a
particle moving once around the ellipse
\[
  \frac{x^2}{4}+\frac{y^2}{9}=1
\]
in the \textbf{clockwise} direction.

\item Consider the vector field
\[
  \vec F(x,y)=\bigl(ye^x+2x,\,e^x+2y\bigr)
\]
and the curve $C$ along the parabola $y=x^2+1$, oriented from $(0,1)$
to $(1,2)$.

\begin{enumerate}[label=(\alph*),leftmargin=1.8em,topsep=0.4em,itemsep=0.2em]
  \item Determine whether $\vec F$ is conservative.
  \item If $\vec F$ is conservative, find a potential
        function $f$. Then use the Fundamental Theorem of Line Integrals
        to evaluate $\displaystyle\int_C\vec F\cdot\dd\vec r$.
\end{enumerate}

\end{enumerate}

\end{document}
